\section{The Paper's Additional Security Claims}\label{sec:claims}

Beyond the proven unforgeability, the paper's Section~V-C lists four further security properties. Two are
testable and we tested them.

\subsection{``Resist replay attacks''}
The paper's argument, in full: ``The signature algorithm of our CBAS scheme is probabilistic, so the
different random values will be used every time the algorithm is run, and the different signature will
be obtained even the same message takes as the input.''~\cite{qiao2023}

Signing is indeed randomised. It does not follow that replays are resisted: a replay attack does not sign
anything, it re-sends an old valid message and signature, and that pair still verifies. On both backends
an old signature verified again, and it also verified inside a new aggregate, even when the same signature
appeared twice in one batch (2/2 each). The paper sets $\Delta$ once in \textsf{Setup} and never says it
changes; if a deployment rotated $\Delta$ per epoch, old signatures would stop verifying in the next epoch,
but replays within an epoch would still succeed. Replay protection needs freshness that the verifier
checks, such as a sequence number or timestamp in the message, or rejection of repeated $T_i$. Replay is a
protocol-level property; randomised signing does not provide it. \textbf{The claim does not hold as
argued.}

\subsection{``Perfect forward security''}
The paper's argument, in full: ``In our scheme, no one can obtain the random value that used to created
the signature, hence, our proposal provides perfect forward security.''~\cite{qiao2023}

The paper gives no definition and no proof. For signatures, forward security has a standard
meaning~\cite{bellare1999}: secret keys evolve through time periods, so a key stolen later cannot sign for
an earlier period. The repaired scheme has no time periods and no key update; $sk_i$ and $\mathit{Cert}_i$
are fixed for life. We tested the consequence: keys recovered while the system was in epoch~1 (through the
nonce reuse of \cref{sec:nonce}) signed a back-dated epoch-0 reading that verified, also inside an
aggregate with a genuine epoch-0 signature (2/2). The argument's premise also fails: ``no one can obtain the
random value'' assumes a perfect generator, and with a repeated nonce $t$ is recovered exactly. Nonce
secrecy is needed for unforgeability; it does not give forward security. The paper itself cites a scheme
titled as forward-secure and certificate-based~\cite{li2015}, so the notion was available to the authors.
\textbf{The claim is not supported}; making it true would need a key-evolving construction, not a change
of wording.
